% PeerJ Computer Science Submission -- Cortex-Nexus
% Using wlpeerj class (official Overleaf template)
\documentclass[lineno]{wlpeerj}

\usepackage{multirow}
\usepackage{float}
\usepackage{siunitx}
\sisetup{
  round-mode=places,
  round-precision=3,
  detect-weight=true,
  detect-family=true,
}

\keywords{emotional prompt engineering, large language models, domain-specific optimization, code generation, LLM-as-Judge, self-refinement}

\title{Cortex-Nexus: Domain-Specific Emotional Prompt Engineering for Large Language Models --- A Systematic Empirical Study Across Philosophical, Technical, and Code Generation Tasks}

\author[1]{Maximiliano Rodrigo Speranza\thanks{Corresponding author: maximiliano.speranza@gmail.com}}
\affil[1]{Independent Researcher, Buenos Aires, Argentina}

\begin{abstract}
Recent work on emotional prompt engineering (EmotionPrompt) has demonstrated that embedding affective cues in system prompts can improve large language model (LLM) performance on benchmarks. However, existing studies apply uniform emotional stimuli without considering domain-specific cognitive demands. We present Cortex-Nexus, a systematic empirical program comprising eight experimental phases (V3C--V6B, $n = 912$ total evaluation cycles) that investigates domain-specific emotional optimization across three qualitatively distinct task domains: philosophical reasoning, technical analysis, and Python code generation. Using an automated adversarial evaluation loop with LLM-as-Judge scoring, we demonstrate that (1) each domain has a distinct optimal emotional configuration --- curiosity for philosophical tasks ($\Delta = +0.192$, $p = 0.022$), concentration for technical tasks ($\Delta = +0.178$, $p < 0.05$), and technical mastery for code generation ($\Delta = +0.357$, $p = 0.003$, $d = 0.528$); (2) the optimal emotion for one domain can be suboptimal or even detrimental in another; (3) combining multiple emotional axes produces cognitive interference that degrades performance ($\Delta = +0.131$, $p = 0.183$, non-significant); and (4) subtle emotional intensity outperforms extreme intensity in code generation ($\Delta_{\text{low}} = +0.404$ vs.\ $\Delta_{\text{extreme}} = +0.163$). We formalize these findings as the \textit{Task-Emotion Alignment Hypothesis}, proposing that emotional prompt engineering must be calibrated to the specific cognitive demands of each task type. All data and code are publicly available.
\end{abstract}

\begin{document}

\flushbottom
\maketitle
\thispagestyle{empty}

\section{Introduction}

Large Language Models (LLMs) have demonstrated remarkable capabilities across diverse tasks, from creative writing to complex code generation \citep{brown2020language, chen2021evaluating}. A growing body of research explores how the framing and structure of prompts influence model performance \citep{wei2022chain, sahoo2024systematic, white2023prompt}. Within this landscape, \citet{coda2023emotionprompt} introduced EmotionPrompt, demonstrating that embedding emotional cues in prompts can yield measurable improvements on standard benchmarks.

However, a critical limitation persists in the emotional prompt engineering literature: existing studies treat emotional stimuli as domain-agnostic --- applying uniform affective cues regardless of whether the target task demands creative exploration, analytical rigor, or algorithmic precision. This one-size-fits-all approach ignores a fundamental insight from cognitive psychology: different cognitive tasks engage different emotional-cognitive pathways \citep{kahneman2011thinking}.

We hypothesize that the optimal emotional configuration for an LLM is not universal but \textit{domain-specific} --- that each task type has an emotional ``sweet spot'' aligned with its unique cognitive demands. We term this the \textbf{Task-Emotion Alignment Hypothesis}.

To test this hypothesis, we developed Cortex-Nexus, an automated experimental platform that conducts controlled comparisons between emotionally-prompted LLMs and neutral baselines across three qualitatively distinct domains:

\begin{enumerate}
    \item \textbf{Philosophical reasoning}: tasks demanding creative exploration, conceptual synthesis, and speculative depth.
    \item \textbf{Technical analysis}: tasks demanding precision, completeness, and structured clarity.
    \item \textbf{Code generation}: tasks demanding algorithmic correctness, efficiency, robustness, and idiomatic expression.
\end{enumerate}

Over eight experimental phases spanning approximately 912 evaluation cycles, Cortex-Nexus systematically explored the interaction between emotional states and task domains, producing what we believe is the most comprehensive empirical dataset on domain-specific emotional prompt engineering to date.

\subsection{Contributions}

This work makes the following contributions:

\begin{enumerate}
    \item We present the \textbf{Task-Emotion Alignment Hypothesis}, formalizing the intuition that different task types require different emotional configurations.
    \item We provide \textbf{systematic empirical evidence} across three domains (philosophical, technical, code) demonstrating that optimal emotions are domain-specific and non-transferable.
    \item We introduce a \textbf{self-refinement pipeline} (Pass 1: emotional generation $\rightarrow$ Pass 2: neutral critique $\rightarrow$ Pass 3: neutral refinement) that amplifies the emotional effect in code generation from $\Delta = +0.050$ to $\Delta = +0.357$.
    \item We document \textbf{cognitive interference}: combining multiple emotional axes degrades rather than enhances performance, establishing practical upper bounds for emotional prompt complexity.
    \item We discover a \textbf{subtlety effect}: low-intensity emotional prompts consistently outperform extreme intensity, contradicting the naive assumption that ``more emotion means better performance.''
\end{enumerate}

\section{Related Work}

\subsection{Prompt Engineering}

The field of prompt engineering has evolved rapidly from simple instruction-following to sophisticated techniques such as chain-of-thought reasoning \citep{wei2022chain}, self-consistency, and tree-of-thought approaches \citep{sahoo2024systematic}. \citet{white2023prompt} cataloged prompt patterns that enhance LLM performance through structural modifications. Our work extends this literature by exploring \textit{affective} rather than structural prompt modifications.

\subsection{Emotional Prompt Engineering}

\citet{coda2023emotionprompt} demonstrated that emotional stimuli embedded in prompts (e.g., ``This is very important to my career'') improve LLM performance on benchmarks by 8--115\%. \citet{li2024more} extended this work, showing that LLMs can both understand and be enhanced by emotional stimuli. However, both studies applied uniform emotional stimuli without domain-specific calibration. Our work addresses this gap by systematically varying emotions across domains and measuring domain-specific interactions.

\subsection{LLM-as-Judge Evaluation}

\citet{zheng2023judging} established the LLM-as-Judge paradigm, demonstrating that strong LLMs can serve as reliable evaluators when properly calibrated. We adopt this approach with randomized presentation order to mitigate positional bias, using Qwen-2.5-72B as our evaluation judge across all phases after V3.

\subsection{Self-Refinement}

\citet{madaan2023self} introduced Self-Refine, showing that iterative self-feedback loops improve LLM output quality. Our work integrates self-refinement with emotional prompting, applying emotion exclusively to the initial generation pass while keeping critique and refinement passes neutral --- a novel combination that amplifies emotional effects without contaminating the refinement process.

\section{Methodology}

\subsection{Experimental Architecture}

Cortex-Nexus operates as a fully automated experimental loop (Figure~\ref{fig:architecture}) with three core components:

\begin{enumerate}
    \item \textbf{Question Generator (Interrogator)}: Generates novel task prompts using a lightweight model (Gemma-3-4B) with deduplication via MD5 hashing to prevent question repetition (a critical fix implemented after V2, which suffered 82.6\% question repetition).
    \item \textbf{Generative Agent (Ente)}: Produces responses using Llama-3.3-70B-Instruct via NVIDIA NIM or OpenRouter APIs. In each cycle, the Ente generates two responses: one with the experimental emotional prompt and one with a neutral control prompt.
    \item \textbf{Evaluation Judge (Juez)}: Evaluates both responses blindly using Qwen-2.5-72B-Instruct with randomized presentation order. Scores are assigned on a 1--10 scale across domain-specific quality dimensions.
\end{enumerate}

\begin{figure}[ht]
\centering
\fbox{\parbox{0.85\linewidth}{\centering
\textbf{Cortex-Nexus Experimental Architecture}\\[6pt]
\small
Question Generator (Gemma-3-4B) $\rightarrow$ MD5 Dedup Buffer\\[3pt]
$\downarrow$\\[3pt]
Ente (Llama-3.3-70B) $\times$ 2: [Emotional] + [Neutral]\\[3pt]
$\downarrow$ \textit{(optional: Self-Refinement Pipeline)}\\[3pt]
\textit{Pass 1}: Emotional Generation $\rightarrow$ \textit{Pass 2}: Neutral Critique $\rightarrow$ \textit{Pass 3}: Neutral Refinement\\[3pt]
$\downarrow$\\[3pt]
Judge (Qwen-2.5-72B) --- Blind, randomized order\\[3pt]
$\downarrow$\\[3pt]
$\Delta = \text{Score}_{\text{exp}} - \text{Score}_{\text{ctl}}$ $\rightarrow$ JSONL persistence
}}
\caption{Architecture of the Cortex-Nexus automated evaluation loop. The self-refinement pipeline was introduced in V3C and used in all code generation experiments from V5C onward.}
\label{fig:architecture}
\end{figure}

\subsection{Evaluation Dimensions}

\subsubsection{Philosophical Domain}
Three dimensions: Depth, Originality, Coherence. Quality score computed as unweighted average.

\subsubsection{Technical Domain}
Three dimensions: Precision, Completeness, Clarity. Quality score computed as unweighted average.

\subsubsection{Code Generation Domain}
Six dimensions with asymmetric weights reflecting practical importance: Correctness (25\%), Efficiency (20\%), Robustness (15\%), Readability (15\%), Simplicity (15\%), and Idiomaticity (10\%).

\subsection{Self-Refinement Pipeline}

Starting from V3C, code generation experiments employed a three-pass self-refinement pipeline:

\begin{enumerate}
    \item \textbf{Pass~1 (Emotional)}: The Ente generates initial code under the experimental emotional prompt.
    \item \textbf{Pass~2 (Neutral Critique)}: A neutral system prompt directs the same model to identify bugs, edge-case failures, complexity inefficiencies, and PEP-8 violations.
    \item \textbf{Pass~3 (Neutral Refinement)}: The model rewrites the code incorporating all identified corrections.
\end{enumerate}

Critically, emotion is applied \textit{only} to Pass~1. Passes 2 and 3 use neutral prompts for both experimental and control conditions. This design isolates the creative/generative effect of emotion from the analytical refinement process.

\subsection{Experimental Phases}

Table~\ref{tab:phases} summarizes the eight experimental phases conducted over the course of the research program.

\begin{table}[H]
\centering
\caption{Summary of experimental phases. Each phase tested specific hypotheses about emotional prompt effects.}
\label{tab:phases}
\small
\begin{tabular}{@{}llllr@{}}
\toprule
\textbf{Phase} & \textbf{Domain} & \textbf{Focus} & \textbf{Key Variable} & \textbf{$n$} \\
\midrule
V3C & Code & Self-refinement pilot & Perfectionism vs.\ Minimalism & 128 \\
V4A & Technical & Curiosity ladder & Intensity (0.20--0.95) & 122 \\
V4B & Code & Perfectionism test & Perfectionism vs.\ Neutral & 91 \\
V5A & Technical & Cross-domain transfer & Phil.\ optimum in technical & 123 \\
V5B & Technical & Bidimensional test & Concentration vs.\ optimum & 98 \\
V5C & Code & Mastery discovery & Mastery vs.\ Precision & 116 \\
V6A & Code & Intensity ladder & Low / Med / High / Extreme & 115 \\
V6B & Code & Synergy test & Mastery $\times$ Precision & 119 \\
\midrule
 & & & \textbf{Total} & \textbf{912} \\
\bottomrule
\end{tabular}
\end{table}

\subsection{Statistical Analysis}

All comparisons use the one-sample $t$-test on the per-cycle delta ($\Delta = \text{Score}_{\text{exp}} - \text{Score}_{\text{ctl}}$) against $H_0: \mu_\Delta = 0$. Effect size is reported as Cohen's $d$ ($d = \bar{\Delta} / s_{\Delta}$). Significance threshold is $\alpha = 0.05$. We report exact $p$-values throughout and do not apply multiple comparison corrections, as each phase tests a distinct, pre-specified hypothesis rather than fishing for significance across many comparisons.

\subsection{Computing Infrastructure}

All experiments were conducted on a consumer-grade workstation running Ubuntu 24.04 LTS. The generative agent (Llama-3.3-70B-Instruct) and evaluation judge (Qwen-2.5-72B-Instruct) were accessed exclusively through cloud APIs: NVIDIA NIM (\texttt{integrate.api.nvidia.com}) and OpenRouter (\texttt{openrouter.ai}). The question generator (Gemma-3-4B) was run locally. The experimental engine was implemented in Python 3.12, with data persisted in JSONL format. Monitoring dashboards were built with Streamlit. The codebase is available in the associated repository.

\section{Results}

\subsection{Exploratory Phases (V3C--V5B)}

Table~\ref{tab:exploratory} summarizes the results of the five exploratory phases that preceded the confirmatory experiments. These phases served to narrow the emotional configuration space and identify promising candidates.

\begin{table}[H]
\centering
\caption{Exploratory phase results (V3C--V5B). Phases V3C and V4B explored code generation; V4A, V5A, and V5B explored the technical domain. Bold indicates statistically significant results ($p < 0.05$).}
\label{tab:exploratory}
\small
\begin{tabular}{@{}llrlrrr@{}}
\toprule
\textbf{Phase} & \textbf{Condition} & \textbf{$n$} & $\bar{\Delta}$ & $s_\Delta$ & \textbf{Direction} \\
\midrule
\multicolumn{6}{l}{\textit{V3C: Self-Refinement Pilot (Code)}} \\
& Perfectionism (0.85) & 61 & $+0.257$ & 1.512 & Positive \\
& Minimalism (0.90) & 67 & $-0.082$ & 1.857 & Negative \\
\midrule
\multicolumn{6}{l}{\textit{V4A: Curiosity Intensity Ladder (Technical)}} \\
& Low (0.20) & 14 & $-0.179$ & 1.185 & Negative \\
& Medium (0.50) & 25 & $+0.160$ & 0.530 & Positive \\
& High (0.75) & 37 & $-0.157$ & 0.843 & Negative \\
& Extreme (0.95) & 40 & $+0.028$ & 0.842 & Neutral \\
\midrule
\multicolumn{6}{l}{\textit{V4B: Perfectionism Validation (Code)}} \\
& Perfectionism (0.85) & 43 & $+0.036$ & -- & Marginal \\
& Neutral baseline & 48 & $-0.059$ & -- & Negative \\
\midrule
\multicolumn{6}{l}{\textit{V5A: Cross-Domain Transfer (Technical)}} \\
& Phil.\ Optimum (C+F) & 94 & $-0.065$ & 0.841 & Negative \\
& Neutral control & 24 & $+0.175$ & 0.705 & Positive \\
\midrule
\multicolumn{6}{l}{\textit{V5B: Bidimensional Comparison (Technical)}} \\
& Phil.\ Optimum (C+F) & 45 & $+0.138$ & 0.786 & Positive \\
& Concentration (0.50) & 45 & $-0.036$ & 0.781 & Neutral \\
\bottomrule
\end{tabular}
\end{table}

Key findings from the exploratory phases: (a)~self-refinement amplified emotional effects relative to zero-shot generation (V3C); (b)~the curiosity intensity ladder showed a non-monotonic pattern in the technical domain, with medium intensity (0.50) outperforming extremes (V4A); (c)~the philosophical optimum (Curiosity + Frustration) \textit{underperformed} the neutral control in technical tasks (V5A), providing the first strong evidence against universal emotional optima; and (d)~perfectionism maintained only marginal advantages in code generation (V4B), motivating the search for stronger configurations.

\subsection{Phase V5C: Mastery Discovery (Code Domain)}

This phase introduced ``technical mastery'' --- a novel emotional configuration emphasizing identity as a senior expert --- which produced the strongest effect in the entire program ($n = 116$).

\begin{table}[H]
\centering
\caption{V5C results: discovery of technical mastery as the optimal code generation emotion.}
\label{tab:v5c}
\begin{tabular}{@{}lrrrrrr@{}}
\toprule
\textbf{Condition} & \textbf{$n$} & $\bar{\Delta}$ & $s_\Delta$ & $t$ & $p$ & $d$ \\
\midrule
\textbf{Mastery (0.85)} & \textbf{37} & $\mathbf{+0.357}$ & \textbf{0.676} & \textbf{3.210} & $\mathbf{0.003}$ & $\mathbf{0.528}$ \\
Precision (0.85) & 49 & $+0.150$ & 0.516 & 2.035 & 0.047 & 0.291 \\
Perfectionism (0.85) & 30 & $-0.158$ & 0.518 & $-1.673$ & 0.105 & $-0.305$ \\
\bottomrule
\end{tabular}
\end{table}

Technical mastery ($d = 0.528$, a medium effect by Cohen's standards) was the breakthrough finding. The mastery prompt frames the model as a ``senior engineer with deep expertise,'' activating what we interpret as a confidence-competence cognitive mode. Notably, perfectionism --- previously the best-performing code emotion --- reversed sign to negative when compared against mastery.

\subsection{Phase V6A: Mastery Intensity Ladder (Code Domain)}

Having identified mastery as the optimal code emotion, V6A explored the dose-response curve across four intensity levels ($n = 115$).

\begin{table}[H]
\centering
\caption{V6A results: mastery intensity ladder. Lower intensity paradoxically produces larger effects.}
\label{tab:v6a}
\begin{tabular}{@{}lrrrrrr@{}}
\toprule
\textbf{Intensity} & \textbf{$n$} & $\bar{\Delta}$ & $s_\Delta$ & $t$ & $p$ & $d$ \\
\midrule
\textbf{Low (0.30)} & \textbf{27} & $\mathbf{+0.404}$ & 0.903 & 2.323 & 0.028 & 0.447 \\
Medium (0.55) & 29 & $+0.371$ & 0.726 & 2.751 & 0.010 & 0.511 \\
High (0.80) & 29 & $+0.205$ & \textbf{0.360} & \textbf{3.068} & \textbf{0.005} & \textbf{0.570} \\
Extreme (0.95) & 30 & $+0.163$ & 0.504 & 1.776 & 0.086 & 0.324 \\
\bottomrule
\end{tabular}
\end{table}

A striking pattern emerged: while all intensity levels produced positive deltas, the \textit{lowest} intensity (0.30) produced the \textit{largest} raw effect ($\Delta = +0.404$), and the extreme intensity (0.95) not only showed the smallest effect but was the \textit{only} level that failed to reach statistical significance ($p = 0.086$). We term this the ``subtlety effect'': emotional prompt engineering is more effective when applied gently.

The high intensity level (0.80) achieved the best Cohen's $d$ ($d = 0.570$) due to its notably low variance ($s = 0.360$), indicating the most consistent (though not the largest) performance improvement.

\subsection{Phase V6B: Bidimensional Synergy Test (Code Domain)}

The final phase tested whether combining mastery with algorithmic precision would produce synergistic or interfering effects ($n = 119$).

\begin{table}[H]
\centering
\caption{V6B results: combining emotions produces cognitive interference.}
\label{tab:v6b}
\begin{tabular}{@{}lrrrrrr@{}}
\toprule
\textbf{Condition} & \textbf{$n$} & $\bar{\Delta}$ & $s_\Delta$ & $t$ & $p$ & $d$ \\
\midrule
Mastery (0.85) & 39 & $+0.308$ & 0.671 & 2.863 & 0.007 & 0.459 \\
Precision (0.85) & 40 & $+0.179$ & 0.463 & 2.444 & 0.019 & 0.386 \\
\textbf{Combined (M+P)} & \textbf{40} & $+0.131$ & 0.612 & 1.356 & \textbf{0.183} & 0.214 \\
\bottomrule
\end{tabular}
\end{table}

The combined condition produced the \textit{weakest} effect and was the only condition that failed to reach statistical significance ($p = 0.183$). This demonstrates \textbf{cognitive interference}: when the system prompt contains multiple emotional directives, the model appears unable to fully activate either mode, resulting in a diluted effect.

\subsection{Cross-Domain Summary}

Table~\ref{tab:alignment} consolidates the optimal emotional configuration for each domain.

\begin{table}[H]
\centering
\caption{The Task-Emotion Alignment map: optimal emotions by domain.}
\label{tab:alignment}
\begin{tabular}{@{}llrrrll@{}}
\toprule
\textbf{Domain} & \textbf{Optimal Emotion} & $\bar{\Delta}$ & $p$ & $d$ & \textbf{Cognitive Mode} \\
\midrule
Philosophical & Curiosity (0.95) & $+0.192$ & 0.022 & -- & Exploration \\
Technical & Concentration (0.50) & $+0.178$ & ${<}0.05$ & -- & Focused analysis \\
Code & Mastery (0.30--0.85) & $+0.357$ & 0.003 & 0.528 & Expert confidence \\
\bottomrule
\end{tabular}
\end{table}

\section{Discussion}

\subsection{The Task-Emotion Alignment Hypothesis: Confirmed}

Each domain's optimal emotion aligns with its cognitive demands:
\begin{itemize}
    \item \textbf{Philosophical tasks} reward creative exploration and conceptual novelty. Curiosity activates a divergent-thinking mode that generates more original connections.
    \item \textbf{Technical tasks} reward precision, completeness, and structured reasoning. Concentration activates a convergent-thinking mode that reduces errors.
    \item \textbf{Code generation} rewards correctness, efficiency, and robustness. Technical mastery activates a confidence-competence mode where the model produces more robust and idiomatically correct code.
\end{itemize}

\subsection{The Subtlety Effect}

One of the most counterintuitive findings is that lower emotional intensity consistently outperforms higher intensity in code generation (V6A). We propose two complementary explanations:

\begin{enumerate}
    \item \textbf{Cognitive rigidity}: High-intensity emotional prompts may constrain the model's token-level search space, reducing the diversity of candidate completions and biasing toward stereotypical solutions.
    \item \textbf{Prompt saturation}: As emotional directives consume more of the system prompt, they may crowd out task-relevant information, reducing the model's effective context.
\end{enumerate}

\subsection{Cognitive Interference in Multi-Axis Emotions}

The V6B finding that combining mastery and precision degrades performance has practical implications. We interpret this as a form of cognitive interference analogous to dual-task interference observed in human cognition \citep{kahneman2011thinking}. The model, when asked to simultaneously adopt the identity of a ``technical master'' and ``algorithmic precision expert,'' fails to fully commit to either role.

\subsection{Evolution of the Self-Refinement Pipeline}

The three-pass self-refinement pipeline proved critical for unlocking emotional effects in code generation. Prior to V3C, code generation showed minimal emotional sensitivity ($\Delta = +0.050$ for perfectionism). After introducing self-refinement, the same domain yielded $\Delta = +0.357$ for mastery ($d = 0.528$), a $7\times$ amplification.

\subsection{Limitations}

\begin{enumerate}
    \item \textbf{Single model family}: All experiments used Llama-3.3-70B. The findings may not generalize to other architectures (GPT-4, Claude, Gemini).
    \item \textbf{LLM-as-Judge reliability}: While we used randomized order to mitigate positional bias, LLM judges have known biases. Human evaluation would strengthen the findings.
    \item \textbf{Question distribution}: The question generator's topic distribution may not be fully representative of real-world task distributions.
    \item \textbf{Statistical power}: Some conditions have modest sample sizes ($n = 27$--$49$). Larger-scale replication would increase confidence.
    \item \textbf{Single researcher}: Experimental design, execution, and analysis were conducted by a single researcher, limiting inter-rater reliability.
\end{enumerate}

\subsection{Future Work}

\begin{enumerate}
    \item \textbf{Cross-model replication}: Testing the hypothesis on GPT-4, Claude, and Gemini.
    \item \textbf{Human evaluation}: Parallel human evaluations to validate LLM-as-Judge findings.
    \item \textbf{Temporal dynamics}: Investigating whether the optimal emotion changes with increasing conversational context.
    \item \textbf{Sub-domain mapping}: Extending the taxonomy to finer-grained task categories.
    \item \textbf{Mechanistic interpretation}: Using activation analysis to understand \textit{how} emotional prompts modify internal representations.
\end{enumerate}

\section{Conclusions}

Cortex-Nexus provides systematic empirical evidence that emotional prompt engineering is not a one-size-fits-all technique. Through 912 controlled evaluation cycles across eight experimental phases, we demonstrate that:

\begin{enumerate}
    \item Each task domain has a distinct optimal emotional configuration, confirming the Task-Emotion Alignment Hypothesis.
    \item Technical mastery ($d = 0.528$) produces the strongest and most robust emotional effect in code generation, representing a practically significant improvement.
    \item Subtle emotional prompts outperform extreme ones, establishing the Subtlety Effect as a practical guideline.
    \item Combining multiple emotional axes produces cognitive interference rather than synergy.
    \item Self-refinement pipelines amplify emotional effects by separating creative generation from analytical refinement.
\end{enumerate}

These findings have direct practical applications: by selecting the appropriate emotional prompt for each task type, practitioners can achieve measurable improvements in LLM output quality without any model modification, fine-tuning, or additional computational cost.

\section*{Data and Code Availability}

All experimental data (JSONL format, 912 records across 8 files), source code for the Cortex-Nexus experimental engine, configuration files, and analysis scripts are publicly available at: \texttt{https://github.com/SperanzaMax/Cortex-Nexus}. An archival snapshot with DOI is available at Zenodo: \texttt{https://doi.org/10.5281/zenodo.19752189}.

\section*{Competing Interests}

The author declares that there are no competing interests.

\section*{Author Contributions}

Maximiliano R.\ Speranza conceived and designed the experiments, developed the Cortex-Nexus software platform, performed the experiments, analyzed the data, prepared figures and tables, authored the manuscript, and approved the final draft.

\section*{Funding}

This research received no external funding. The author is an independent researcher with no institutional affiliation or grant support. All computational resources were self-funded. API costs for NVIDIA NIM and OpenRouter were covered by the author's personal accounts.

\section*{Acknowledgments}

The author thanks the open-source AI community for providing access to the Llama-3.3-70B and Qwen-2.5-72B models via NVIDIA NIM and OpenRouter APIs.

\bibliography{peerj}

\end{document}
